\section*{Capítulo \ref{ch:g12:ph-conductivity-titrations} --- Titulações por pH e por condutividade}

\begin{solution}{exo:g12:ph-conductivity-titrations:1}
$V_E = \qty{20.0}{mL}$; $c = 0.100 \times 20.0 / 20.0 = \qty{0.100}{mol/L}$.
\end{solution}

\begin{solution}{exo:g12:ph-conductivity-titrations:2}
Na \omterm{def:g12:ph-conductivity-titrations:ph-metric}{semiequivalência}, \omterm{def:g9:acids-bases-ph:ph}{pH} $= pK_a$: $pK_a = 3.75$, próximo do ácido
metanoico (3.74).
\end{solution}

\begin{solution}{exo:g12:ph-conductivity-titrations:3}
A fenolftaleína: a sua faixa 8.0--10.0 fica dentro do salto, que é
básico. O vermelho de metila (4.2--6.3) mudaria perto da
\omterm{def:g12:ph-conductivity-titrations:ph-metric}{semiequivalência}, cedo demais.
\end{solution}

\begin{solution}{exo:g12:ph-conductivity-titrations:4}
A resposta do eletrodo varia com o tempo e a temperatura: ele é ajustado
para ler corretamente em duas \omterm{def:g12:buffers-predominance:buffer}{soluções-tampão} de \omterm{def:g9:acids-bases-ph:ph}{pH} conhecido.
\end{solution}

\begin{solution}{exo:g12:ph-conductivity-titrations:5}
\ce{H3O+} (349.8) > \ce{OH-} (198.3) > \ce{Cl-} (76.2) > \ce{Na+} (50.3).
\end{solution}

\begin{solution}{exo:g12:ph-conductivity-titrations:6}
$c = 0.050 \times 14.6 / 20.0 = \qty{0.0365}{mol/L}$.
\end{solution}

\begin{solution}{exo:g12:ph-conductivity-titrations:7}
Inclinações: 2.0, 4.5, 17.5 e \qty{2.5}{mL^{-1}}. A maior fica entre
10.0 e \qty{10.2}{mL}: $V_E \approx \qty{10.1}{mL}$.
\end{solution}

\begin{solution}{exo:g12:ph-conductivity-titrations:8}
Antes da \omterm{def:g11:titration:equivalence}{equivalência}, cada \ce{OH-} adicionado elimina um \ce{H3O+}
(muito condutor) e traz um \ce{Na+} (muito menos condutor): a
\omterm{def:g12:ph-conductivity-titrations:conductivity}{condutividade} cai. Depois dela, o \ce{Na+} e o \ce{OH-} se acumulam sem
reagir: ela sobe.
\end{solution}

\begin{solution}{exo:g12:ph-conductivity-titrations:9}
No início, poucos \omterm{def:g9:ions:ion}{íons} (o \omterm{def:g12:ka-and-pka:strong-acid}{ácido fraco} quase não reage com a água):
\omterm{def:g12:ph-conductivity-titrations:conductivity}{condutividade} baixa. Antes da \omterm{def:g11:titration:equivalence}{equivalência}, \omterm{def:g9:ions:ion}{íons} etanoato e sódio se
formam no lugar de \omterm{def:g7:atoms-and-molecules:molecule}{moléculas} sem carga: ela sobe, devagar. Depois da
\omterm{def:g11:titration:equivalence}{equivalência}, o \ce{Na+} e o \ce{OH-} se acumulam: ela sobe mais
depressa. Duas retas que sobem, a segunda mais inclinada.
\end{solution}

\begin{solution}{exo:g12:ph-conductivity-titrations:10}
Cerca de 2.9, 4.76 e 8.7. O ácido etanoico é fraco: só uma pequena
porcentagem dele dá \omterm{def:g12:ka-and-pka:oxonium}{íons oxônio}, por isso o \omterm{def:g9:acids-bases-ph:ph}{pH} começa mais alto que o
1.0 do ácido clorídrico numa concentração parecida.
\end{solution}

\begin{solution}{exo:g12:ph-conductivity-titrations:11}
Para que o volume de \omterm{def:g11:titration:titration}{titulante} adicionado quase não mude o volume total:
as concentrações e, portanto, a \omterm{def:g12:ph-conductivity-titrations:conductivity}{condutividade} variam então ao longo de
retas.
\end{solution}

\begin{solution}{exo:g12:ph-conductivity-titrations:12}
Os \omterm{def:g9:ions:ion}{íons} hidróxido reagem primeiro com o \omterm{def:g12:ka-and-pka:strong-acid}{ácido forte} (\ce{H3O+}), depois
com o fraco. O primeiro salto marca o fim do ácido clorídrico (o seu
\omterm{def:g11:titration:equivalence}{volume de equivalência} o mede), o segundo, o fim do ácido etanoico (o
volume entre os dois saltos o mede).
\end{solution}

\begin{solution}{exo:g12:ph-conductivity-titrations:13}
Início: $(349.8 + 76.2) \times 0.0100 = \qty{4.26}{mS/cm}$. Na
\omterm{def:g11:titration:equivalence}{equivalência}: $[\ce{Na+}] = [\ce{Cl-}] = 1.00/110 = \qty{0.00909}{mol/L}$,
logo $(50.3 + 76.2) \times 0.00909 = \qty{1.15}{mS/cm}$. Os dois como
na figura.
\end{solution}

\begin{solution}{exo:g12:ph-conductivity-titrations:14}
Nenhum em $V_E$: o salto fica no mesmo volume, qualquer que seja o
desvio das leituras. O $pK_a$ lido na \omterm{def:g12:ph-conductivity-titrations:ph-metric}{semiequivalência} fica 0.2 alto
demais.
\end{solution}

\begin{solution}{exo:g12:ph-conductivity-titrations:15}
O \omterm{def:g9:acids-bases-ph:ph}{pH} não muda durante essa \omterm{def:g9:ions:precipitate}{precipitação}, mas os \omterm{def:g9:ions:ion}{íons} sim: antes da
\omterm{def:g11:titration:equivalence}{equivalência}, cada \ce{Ag+} adicionado elimina um \ce{Cl-} e traz um
\ce{NO3-} de \omterm{def:g12:ph-conductivity-titrations:conductivity}{condutividade} parecida, de modo que a \omterm{def:g12:ph-conductivity-titrations:conductivity}{condutividade} fica
quase constante (os \omterm{def:g9:ions:ion}{íons} sódio já estavam lá desde o início); depois da
\omterm{def:g11:titration:equivalence}{equivalência}, o \ce{Ag+} e o \ce{NO3-} se acumulam e ela sobe. Uma reta
horizontal, depois uma reta que sobe.
\end{solution}

\begin{solution}{pb:g12:ph-conductivity-titrations:1}
\textbf{1.} \qty{6}{g} de ácido etanoico por \qty{100}{g} de vinagre.

\textbf{2.} \qty{101}{g} de vinagre contêm $6.0 \times 101/100 =
\qty{6.06}{g}$; $M = \qty{60.0}{g/mol}$: \qty{0.101}{mol} em
\qty{0.1000}{L}, isto é, \qty{1.01}{mol/L}.

\textbf{3.} Sem \omterm{def:g10:concentration-and-dilution:dilution}{diluição}, \qty{10.0}{mL} precisariam de cerca de
\qty{101}{mL} de hidróxido de sódio a \qty{0.100}{mol/L}: mais do que
cabe numa bureta.

\textbf{4.} Uma pipeta volumétrica de \qty{10.0}{mL} e um balão
volumétrico de \qty{100.0}{mL}.

\textbf{5.} \ce{CH3COOH + OH- -> CH3COO- + H2O}.

\textbf{6.} $V_E = \qty{10.1}{mL}$.

\textbf{7.} $c = 0.100 \times 10.1 / 10.0 = \qty{0.101}{mol/L}$.

\textbf{8.} $10 \times 0.101 = \qty{1.01}{mol/L}$.

\textbf{9.} \omterm{def:g9:acids-bases-ph:ph}{pH} 4.76 em \qty{5.05}{mL}: o $pK_a$ do ácido etanoico.

\textbf{10.} Na \omterm{def:g11:titration:equivalence}{equivalência}, a solução contém etanoato de sódio, e o
etanoato é uma \omterm{def:g12:ka-and-pka:strong-acid}{base fraca}.

\textbf{11.} A fenolftaleína, cuja faixa 8.0--10.0 contém o \omterm{def:g9:acids-bases-ph:ph}{pH} na
\omterm{def:g11:titration:equivalence}{equivalência} (cerca de 8.7).

\textbf{12.} \omterm{def:g9:ions:ion}{Íons} sódio e \omterm{def:g9:ions:ion}{íons} etanoato.

\textbf{13.} Cada \ce{OH-} adicionado transforma uma \omterm{def:g7:atoms-and-molecules:molecule}{molécula} quase não
ionizada em \ce{CH3COO-}, com um \ce{Na+}: \omterm{def:g9:ions:ion}{íons} de \omterm{def:g12:ph-conductivity-titrations:conductivity}{condutividade} modesta
são acrescentados aos poucos.

\textbf{14.} Depois da \omterm{def:g11:titration:equivalence}{equivalência}, os \omterm{def:g9:ions:ion}{íons} hidróxido, muito
condutores, se acumulam com os \omterm{def:g9:ions:ion}{íons} sódio.

\textbf{15.} O \omterm{def:g12:ka-and-pka:strong-acid}{ácido fraco} dá muito poucos \omterm{def:g9:ions:ion}{íons} à água.

\textbf{16.} Em $V_E = \qty{10.1}{mL}$, como na \omterm{def:g12:ph-conductivity-titrations:ph-metric}{titulação
potenciométrica}.

\textbf{17.} $1.01 \times 60.0 = \qty{60.6}{g/L}$.

\textbf{18.} \qty{6.06}{g} em \qty{100.0}{mL}.

\textbf{19.} $6.06 \times 100 / 101 = \qty{6.0}{g}$ por \qty{100}{g}.

\textbf{20.} O vinagre tem uma acidez de \qty{6.0}{\percent}: o rótulo está certo.
\end{solution}
